\begin{table}[t]
\centering
\caption{Standardized Effect Sizes}
\label{tab:sde}
\begin{threeparttable}
\begin{tabular}{lcccccc}
\toprule
Outcome & $\hat{\beta}$ & SE & SD($Y$) & SDE & SE(SDE) & Classification \\
\midrule
\multicolumn{7}{l}{\textit{Panel A: Pooled}} \\
All fish density & 0.001 & 0.019 & 0.123 & 0.011 & 0.154 & Small positive \\
Targeted fish density & 0.050 & 0.027 & 0.064 & 0.789 & 0.424 & Large positive \\
Non-targeted density & -0.044 & 0.019 & 0.103 & -0.427 & 0.185 & Large negative \\
Species richness & 5.028 & 1.340 & 5.765 & 0.872 & 0.232 & Large positive \\
\addlinespace
\multicolumn{7}{l}{\textit{Panel B: Heterogeneous (by protection type)}} \\
No-take reserve (IVEE) & 0.016 & 0.006 & 0.065 & 0.238 & 0.099 & Large positive \\
Limited-take area (NAPL) & 0.082 & 0.006 & 0.062 & 1.325 & 0.091 & Large positive \\
\bottomrule
\end{tabular}
\begin{tablenotes}[flushleft]
\footnotesize
\item \textit{Notes:} \textbf{Country:} United States (California). \textbf{Research question:} Does marine protected area designation increase kelp forest fish abundance, and do strict no-take reserves produce larger recovery than limited-take areas? \textbf{Policy mechanism:} The Central Coast Marine Life Protection Act (MLPA, September 2007) designated marine protected areas that restrict or prohibit fishing, reducing fishing mortality on resident reef fish and allowing recovery of depleted populations through reduced harvest pressure. \textbf{Outcome definition:} Log fish density (count per 60m$^2$ transect area + 1) from standardized annual SCUBA surveys. \textbf{Treatment:} Binary---site inside vs.\ outside MPA boundary after September 2007 designation. \textbf{Data:} SBC LTER KFCD annual kelp forest surveys (knb-lter-sbc.17), 2001--2025, 9 mainland reef sites. \textbf{Method:} Two-way fixed effects DiD with site and year FE, standard errors clustered at site level; wild cluster bootstrap and randomization inference for few-cluster robustness. \textbf{Sample:} 9 mainland Santa Barbara Channel kelp forest reefs; Channel Islands excluded (always treated). SDE $= \hat{\beta} / \text{SD}(Y)$ where SD($Y$) is the pre-treatment standard deviation. Classification refers to magnitude, not statistical significance: Large ($|$SDE$| > 0.15$), Moderate ($0.05$--$0.15$), Small ($0.005$--$0.05$), Null ($< 0.005$).
\end{tablenotes}
\end{threeparttable}
\end{table}
